\subsection{Examples}

\begin{enumerate}
\def\labelenumi{\alph{enumi})}
\tightlist
\item
  Let X be a topological space. Let \(\mathcal{R}(\mathrm{X})\) , the vector space of numerical (finite) functions on X be assigned the topology of compact convergence (GT, X, §1, No.~3): this is the coarsest topology for which the restriction mappings \(\mathcal{R}(\mathrm{X}) \to \mathcal{R}(\mathrm{H})\) are continuous (where H runs through the family of compact subsets of X and where \(\mathcal{R}(\mathrm{H})\) is assigned the topology of uniform convergence). Cor. 3 of III, p.~4 shows that a subset A of \(\mathcal{R}(\mathrm{X})\) is bounded if and only if, for every compact subset H of X, the set of restrictions to H of functions belonging to A is uniformly bounded.
\end{enumerate}

* b) (Spaces of infinitely differentiable functions.) Let \(n \geqslant 1\) be an integer. For every open set U in \(\mathbf{R}^n\) , let \(\mathcal{C}^\infty(\mathrm{U})\) denote the vector space of infinitely differentiable functions on U (VAR, R, 2.3). Let \(f\) be in \(\mathcal{C}^\infty(\mathrm{U})\) . For every multi-index \(\alpha = (\alpha_1, \dots, \alpha_n)\) in \(\mathbf{N}^n\) , \(\partial^\alpha f\) denotes the partial derivative \(\partial^{|\alpha|} f / \partial x_1^{\alpha_1} \dots \partial x_n^{\alpha_1}\) ; this is a continuous function on U (VAR, R, 2.3 and 2.4). For every integer \(m \geqslant 0\) , and every compact subset H of U, set

\[
p_{m,\mathrm{H}}(f) = \sup_{\substack{|\alpha |\leqslant m\\ x\in \mathrm{H}}}\left|\partial^{\alpha}f(x)\right|.\tag{1}
\]

Then \(p_{m,\mathrm{H}}\) is a semi-norm on \(\mathcal{C}^{\infty}(\mathrm{U})\) .

Let \(\mathcal{C}^{\infty}(\mathrm{U})\) be assigned the topology defined by the semi-norms \(p_{m,\mathrm{H}}\) . This is the coarsest of the topologies for which the mappings \(f\to \partial^{\alpha}f\) from \(\mathcal{C}^{\infty}(\mathrm{U})\) into \(\mathcal{R}(\mathrm{U})\) are continuous, where \(\mathcal{R}(\mathrm{U})\) is assigned the topology of compact convergence. There exists an increasing sequence of compact subsets \((\mathrm{H}_n)_{n\geqslant 0}\) of U whose interiors cover U; the family of semi-norms \(p_{m,\mathrm{H}_n}\) defines the topology of \(\mathcal{C}^{\infty}(\mathrm{U})\) , which then becomes a locally convex metrizable space. The space \(\mathcal{C}^{\infty}(\mathrm{U})\) is complete; in other words, it is a Fréchet space (II, p.~24): in fact, let \((f_k)\) be a Cauchy sequence in \(\mathcal{C}^{\infty}(\mathrm{U})\) ; for every \(\alpha \in \mathbf{N}^n\) , the sequence \((\partial^{\alpha}f_k)\) converges in the complete space \(\mathcal{R}(\mathrm{U})\) (TG, X, §1, No.~5, th. 1) to a continuous function \(g_{\alpha}\) . By induction on \(|\alpha|\) , we deduce from th. 1 of FVR, II, p.~2 that \(g_{\alpha} = \partial^{\alpha}g_{0}\) for every \(\alpha \in \mathbf{N}^n\) . In other words, the sequence \((f_k)\) converges to \(g_0\) in \(\mathcal{C}^{\infty}(\mathrm{U})\) .

Let A be a subset of \(\mathcal{C}^{\infty}(\mathrm{U})\) . In order that A be bounded, it is necessary and sufficient that the number \(\sup_{f\in \mathbf{A}}p_{m,\mathrm{H}}(f)\) is finite for every integer \(m\geqslant 0\) and for every

  compact subset H of U; this condition means that for every \(\alpha\in \mathbf{N}^{n}\) , the set of functions \(\partial^{\alpha}f|H\) for \(f\in A\) is uniformly bounded for every compact \(H\subset U\) .

Let \(\mathrm{H} \subset \mathrm{U}\) be compact. We denote by \(\mathcal{C}_{\mathrm{H}}^{\infty}(\mathrm{U})\) the subspace of \(\mathcal{C}^{\infty}(\mathrm{U})\) consisting of those functions whose support lies in \(\mathrm{H}\) . The space \(\mathcal{C}_{\circ}^{\infty}(\mathrm{U})\) of infinitely differentiable functions with compact support in \(\mathrm{U}\) is the directed increasing union of subspaces \(\mathcal{C}_{\mathrm{H}}^{\infty}(\mathrm{U})\) where \(\mathrm{H}\) runs through the family of compact subsets of \(\mathrm{U}\) . Each space \(\mathcal{C}_{\mathrm{H}}^{\infty}(\mathrm{U})\) will be assigned the topology induced by that of \(\mathcal{C}^{\infty}(\mathrm{U})\) , and \(\mathcal{C}_{\circ}^{\infty}(\mathrm{U})\) with the corresponding inductive limit topology. If the sets \(\mathrm{H}_n\) are such that their interiors form a cover of \(\mathrm{U}\) , then the space \(\mathcal{C}_{\circ}^{\infty}(\mathrm{U})\) is the strict inductive limit of the Fréchet spaces \(\mathcal{C}_{\mathrm{H}_n}^\infty (\mathrm{U})\) ; it is therefore complete (II, p.~32, prop. 9) and every bounded subset of \(\mathcal{C}_{\circ}^{\infty}(\mathrm{U})\) is contained in one of the subspaces \(\mathcal{C}_{\mathrm{H}_n}^\infty (\mathrm{U})\) (III, p.~5, prop. 6). *

\begin{enumerate}
\def\labelenumi{\alph{enumi})}
\setcounter{enumi}{2}
\tightlist
\item
  (Gevrey's spaces.) Let I be a compact interval in \(\mathbf{R}\) . For every integer \(n \geqslant 0\) , \(\mathrm{D}^n f\) denotes the \(n\) th derivative of a numerical function \(f\) defined on I (whenever this derivative exists). Let \(s \geqslant 1\) and \(\mathbf{M} \geqslant 0\) be two real numbers. Let \(\mathcal{G}_{s,\mathbf{M}}(\mathrm{I})\) denote the vector space of those infinitely differentiable functions \(f\) on I (FVR, I, p.~28) for which the sequence \((|\mathrm{D}^n f| / \mathrm{M}^n (n!)^s)_{n \geqslant 0}\) is bounded in the space \(\mathcal{C}(\mathrm{I})\) of all continuous functions on I (with the topology of uniform convergence). The space \(\mathcal{G}_{s,\mathbf{M}}(\mathrm{I})\) is a Banach space with the norm
\end{enumerate}

\[
\| f \| _ {s, \mathbf {M}} = \sup _ {n \geqslant 0, x \in \mathrm{I}} | \mathrm{D} ^ {n} f (x) | / \mathrm{M} ^ {n} (n!) ^ {s}.
\]

For \(\mathbf{M} \leqslant \mathbf{M}'\) , we have \(\mathcal{G}_{s,\mathbf{M}}(\mathrm{I}) \subset \mathcal{G}_{s,\mathbf{M}'}(\mathrm{I})\) , and

\[
\| f \| _ {s, \mathbf {M} ^ {\prime}} \leqslant \| f \| _ {s, \mathbf {M}}
\]

for every \(f \in \mathcal{G}_{s,\mathsf{M}}(\mathrm{I})\) . Let \(\mathcal{G}_s(\mathrm{I})\) denote the union of the spaces \(\mathcal{G}_{s,\mathsf{M}}(\mathrm{I})\) and endow it with the inductive limit topology of the topologies of \(\mathcal{G}_{s,\mathsf{M}}(\mathrm{I})\) .

Let \(M < M'\) and let B be the unit ball (closed) in \(\mathcal{G}_{s,M}(I)\) . We will prove that B is a compact subset of the Banach space \(\mathcal{G}_{s,M'}(I)\) . It is clear that B is closed in \(\mathcal{G}_{s,M'}(I)\) and so it is enough to prove that B is precompact in \(\mathcal{G}_{s,M'}(I)\) . Let \(\varepsilon > 0\) and let N be a positive integer such that \((M/M')^N \leqslant \varepsilon/2\) . Let k be a positive integer; the set of all functions \(D^{k+1}f\) , as f ranges over B, is bounded in \(\mathcal{C}(I)\) , hence the set of all functions \(D^k f\) , as f ranges over B, is relatively compact in \(\mathcal{C}(I)\) : this follows from the theorem of finite increments (FVR, I, p.~23, cor. 1) and from Ascoli's theorem (GT, X, § 2, No.~5). We define a norm on \(\mathcal{G}_{s,M}(I)\) by

\[
q(f) = \sup_{\substack{0\leqslant n\leqslant \mathbf{N}\\ x\in \mathrm{I}}}\big|\mathrm{D}^{n}f(x)\big| / \mathrm{M}^{\prime n}(n!)^{s}  .
\]

The above argument shows that B is precompact for the topology associated with the norm q; in other words, there exists a finite subset C of B such that for every \(f \in B\) , there exists \(g \in C\) such that \(q(f - g) \leqslant \varepsilon\) . Finally, for every n \textgreater{} N, we have

\[
\left| \mathrm{D} ^ {n} f (x) - \mathrm{D} ^ {n} g (x) \right| / \mathrm{M} ^ {\prime n} (n!) ^ {s} \leqslant 2 (\mathrm{M} / \mathrm{M} ^ {\prime}) ^ {n} \leqslant \varepsilon ,
\]

from which we get \(\| f - g\| \leqslant \varepsilon\) . This proves that B is precompact in \(\mathcal{G}_{s,M'}(I)\) .

The space \(\mathcal{G}_s(\mathrm{I})\) is the inductive limit of the spaces \(\mathcal{G}_{s,k}(\mathrm{I})\) as \(k\) ranges over \(\mathbf{N}\) ; by prop. 7 (III, p.~6) every bounded subset of \(\mathcal{G}_s(\mathrm{I})\) is contained in one of the spaces \(\mathcal{G}_{s,k}(\mathrm{I})\) and is relatively compact in this space.

* \(d\) ) (Spaces of holomorphic functions.) Let \(n \geqslant 1\) be an integer. For every open subset U of \(\mathbf{C}^n\) , \(\mathcal{H}(\mathrm{U})\) denotes the space of functions holomorphic in U, and is assigned the topology of compact convergence in U. For every compact subset L of \(\mathbf{C}^n\) , \(\mathcal{H}(\mathrm{L})\) denotes the space of germs of holomorphic functions in a neighbourhood of L; we endow this space with the finest locally convex topology for which the canonical mappings \(\pi_{\mathrm{U}}:\mathcal{H}(\mathrm{U})\to\mathcal{H}(\mathrm{L})\) are continuous, where U ranges over the set of open neighbourhoods of L.

For every integer \(m \geqslant 1\) , let \(U_{m}\) be the set of points of \(C^{n}\) which are at a distance \textless{} 1/m from L. It can be shown that the canonical mapping \(\pi_{U_{m}}\) from \(\mathcal{H}(U_{m})\) into \(\mathcal{H}(L)\) is injective, and that the restriction mapping from \(\mathcal{H}(U_{m})\) into \(\mathcal{H}(U_{p})\) is compact for \(p \geqslant m\) . We can then apply prop. 7 (III, p.~6). Let A be a bounded subset of \(\mathcal{H}(L)\) ; then there exists an integer \(m \geqslant 1\) such that A consists of germs of functions in a neighbourhood of L, belonging to a bounded set B in \(\mathcal{H}(U_{m})\) . Moreover, a mapping \(\phi\) from \(\mathcal{H}(L)\) into a topological space T is continuous if and only if the mapping \(\phi \circ \pi_{U}\) from \(\mathcal{H}(U)\) into T is continuous for every open neighbourhood U of L. *

